\section*{Chapitre \ref{ch:pl:python-at-scale} --- Python à grande échelle}

\begin{solution}{exo:pl:python-at-scale:1}
$960 \times 2\times10^8$~ns $= 192$~s, soit environ trois minutes~; $145$~ns~: 29~s~; $14$~ns~: 2.7~s.
\end{solution}

\begin{solution}{exo:pl:python-at-scale:2}
Chaque accès à un champ d'un enregistrement projeté crée un objet enregistrement, puis un objet scalaire, avec leur comptage de références
et leur destruction~; la boucle sur les listes lit des entiers qui existent déjà. L'arithmétique est la même~; les objets, non.
\end{solution}

\begin{solution}{exo:pl:python-at-scale:3}
Environ $3 \times 61 = 183$~Mo sous forme d'objets et $3 \times 1.0 = 3$~Mo sous forme de catégorielles.
\end{solution}

\begin{solution}{exo:pl:python-at-scale:4}
Avec $c_0 = 0$ et $c_t = \sum_{s\le t} x_s$ (une somme cumulée), la somme des $w$ dernières valeurs se terminant en $t$ vaut $c_t - c_{t-w}$ (avec
$c_{t-w}$ pris égal à 0 lorsque $t < w$). Une version par blocs reporte les $w-1$ dernières valeurs du bloc précédent (ou, de manière équivalente, les
$w$ dernières sommes cumulées), de sorte que les $w-1$ premières fenêtres d'un bloc soient complètes.
\end{solution}

\begin{solution}{exo:pl:python-at-scale:5}
Deux processus auraient chacun leur propre interpréteur et leur propre verrou, et exécuteraient les deux boucles en parallèle, en environ le temps d'une seule. Avec une
entrée de 64~Mo chacune, sérialisée par tuyau, chacune paierait d'abord environ 0.14~s de \omterm{def:pl:python-at-scale:processes}{coût de sérialisation}, plus que la boucle elle-même~; par mémoire partagée
ou fichier projeté, environ 0.03~s, ou rien si les processus projettent eux-mêmes le fichier.
\end{solution}

\begin{solution}{exo:pl:python-at-scale:6}
$1.5\times10^9 / 16 \approx 94$ millions d'événements, soit environ 0.03~\% d'une journée de 311 milliards de messages~: un flux volumineux ne peut pas être
conservé entier dans un processus~; il doit être lu en flux.
\end{solution}

\begin{solution}{exo:pl:python-at-scale:7}
Les écarts sont de l'ordre de $10^{-13}$ quelle que soit la taille des blocs~: il s'agit d'arrondis, puisque les sommes cumulées de la version par blocs
démarrent à partir de la queue reportée de chaque bloc et non à partir de la première valeur~; elles sont exactes au sens mathématique.
\end{solution}

\begin{solution}{exo:pl:python-at-scale:8}
Une boucle Python pure détient le verrou de l'interpréteur~: huit threads l'exécutent à tour de rôle, en environ le temps d'un seul thread sur la journée
entière, auquel s'ajoute le coût des commutations. Il faut utiliser huit processus, chacun projetant sa part du fichier, ou réécrire la boucle sur des tableaux~; et si la
variable est récursive sur toute la journée, il faut reporter son état d'une partie à l'autre au lieu de le réinitialiser à chaque frontière.
\end{solution}

\begin{solution}{pb:pl:python-at-scale:1}
\begin{enumerate}
\item Elle construit et détruit des objets enregistrement et scalaire, consulte des noms, répartit chaque opération selon les types et compte les références,
pour chacune des dix environ opérations de bytecode d'une itération.
\item Environ 960 (boucle sur les enregistrements), 145 (boucle sur les listes), 14 (expression sur tableaux) et 15 (par blocs).
\item 70 fois plus rapide sur les enregistrements, 11 fois plus rapide sur les listes.
\item Le filtre effectue les multiplications et additions de la récursion dans le même ordre que la boucle.
\item Les boucles dont la forme de l'état change à chaque événement (un carnet d'ordres, une machine à états aux nombreuses branches)~; elles relèvent du code compilé
(chapitre~9).
\item Environ 61, 1, 8 et 16 octets.
\item Cinquante symboles distincts stockés une seule fois, et un code d'un octet par ligne~: \omterm{def:pl:columnar-formats:encodings}{encodage par dictionnaire} en mémoire.
\item Environ 16 octets par événement~: les quantités signées et la sortie, chacune sur 8 octets, plus des variables temporaires de courte durée.
\item Environ 94 millions.
\item Le \omterm{def:pl:python-at-scale:chunked}{pic de mémoire résidente}, soit le plus grand volume de mémoire physique utilisé par le processus à un instant donné~; un job est tué à son pic, non à sa moyenne.
\item La dernière sortie du filtre initialise la récursion du bloc suivant (\cref{prop:pl:python-at-scale:exact}).
\item Des blocs de 100\,000 événements, avec un pic de 2.5~Mo~; des blocs de 10\,000 restent à 10~\% près à 0.31~Mo.
\item Le \omterm{def:pl:python-at-scale:overhead}{surcoût de l'interpréteur} n'est payé qu'une fois par bloc, soit deux mille fois pour deux millions d'événements.
\item Sur une boucle Python, rien (deux fois plus de travail en deux fois plus de temps)~; sur un tri numpy, presque un cœur entier, car le tri libère le verrou
(0.094~s pour deux contre 0.086~s pour un).
\item Environ 0.14~s sérialisé, 0.03~s par mémoire partagée, copie comprise.
\item \textbf{Résultat clé.} L'expression sur tableaux est 70 fois plus rapide qu'une boucle Python sur les enregistrements projetés et 11 fois plus rapide
que celle sur les listes (14 contre 960 et 145 nanosecondes par événement)~; lue en flux par blocs de 10\,000 événements, elle prend à peu près le même temps
(0.025~s contre 0.028~s en bloc unique, 0.023~s au meilleur bloc de 100\,000), avec un pic de 0.31~Mo au lieu de 32~Mo~; un budget de 1.5~Go contient environ 94 millions d'événements en bloc unique, et n'importe quel historique peut être lu en flux.
\item $250 \times 2\times10^8 = 5\times10^{10}$ événements~: environ 13 heures avec la boucle sur les enregistrements, 2 heures sur les listes, 11 minutes avec
les tableaux.
\item Huit processus, chacun projetant sa part des fichiers des jours et les lisant en flux par blocs, chacun renvoyant un petit résultat~; aucune donnée
ne transite par des tuyaux.
\item Cela permettrait aux threads d'exécuter des boucles Python en parallèle, de sorte que la version threadée de l'exercice~8 passerait à l'échelle~; cela ne
changerait ni le \omterm{def:pl:python-at-scale:overhead}{surcoût de l'interpréteur} par événement, ni la mémoire des objets, ni l'intérêt de la lecture en flux.
\item Garder les boucles hors de l'interpréteur et tout l'historique hors de la mémoire.
\end{enumerate}
\end{solution}

\begin{solution}{iq:pl:python-at-scale:1}
Chaque itération s'exécute dans l'interpréteur, en construisant et en vérifiant des objets pour chaque élément~: des centaines de nanosecondes contre quelques-unes
pour l'arithmétique. Il faut l'exprimer avec des opérations sur tableaux (sommes cumulées, filtres, \texttt{searchsorted}), ou compiler la boucle.
\iqlookfor{\omterm{def:pl:python-at-scale:overhead}{Surcoût de l'interpréteur} par élément~; les idiomes sur tableaux.}
\end{solution}

\begin{solution}{iq:pl:python-at-scale:2}
Comme filtre linéaire (un seul appel en code compilé), appliqué aux données par blocs depuis un \omterm{def:pl:a-tick-store-on-open-formats:mmap}{fichier projeté en mémoire}, la dernière sortie de chaque
bloc servant de condition initiale au bloc suivant~; la mémoire reste limitée à un bloc.
\iqlookfor{La récursion comme filtre~; le découpage en blocs avec état reporté.}
\end{solution}

\begin{solution}{iq:pl:python-at-scale:3}
Un verrou qui ne laisse qu'un thread à la fois exécuter du bytecode Python. Il n'a pas d'effet sur les entrées-sorties, ni sur le code compilé qui le libère
(une grande partie de numpy), et on le contourne avec des processus~; une version de CPython sans verrou, disponible dans les versions récentes, le supprime.
\iqlookfor{Concurrence contre parallélisme~; ce qui libère le verrou.}
\end{solution}

\begin{solution}{iq:pl:python-at-scale:4}
Colonnes d'objets (chaînes stockées comme objets Python~: les convertir en catégorielles ou en chaînes Arrow), types trop larges (float64 ou int64 là où 32 ou 8 bits
suffisent), copies conservées par des variables intermédiaires, et question de savoir si l'historique entier doit vraiment tenir en mémoire.
\iqlookfor{Représentations d'abord, puis lecture en flux.}
\end{solution}

\begin{solution}{iq:pl:python-at-scale:5}
On répartit le travail par jours (ou par instruments) dans une file d'attente~; on lance une trentaine de processus de travail, chacun projetant ou parcourant sa partition
depuis la \omterm{def:pl:a-tick-store-on-open-formats:store}{base de ticks} et la lisant en flux par blocs avec un état reporté, et écrivant les résultats dans des fichiers~; on dimensionne les blocs pour que les 30 processus restent
largement dans 128~Go~; on mesure le temps et le pic de mémoire sur un jour avant de tout lancer.
\iqlookfor{Processus plutôt que threads, travail local aux données, budget mémoire mesuré.}
\end{solution}
